% !TeX encoding = UTF-8
\documentclass[variant=guia,cover=institutional,mode=screen]{ua-engineering-v6-article}
% !TeX encoding = UTF-8
\UASetUniversity{Universidad Austral}
\UASetFaculty{Facultad de Ingeniería}
\UASetDepartment{Departamento de Computación}
\UASetCampus{Pilar}
\UASetCareer{Ingeniería Informática}
\UASetCommission{Comisión A}
\UASetProfessor{Equipo docente}
\UASetSemester{Segundo cuatrimestre 2026}
\UASetCourse{Sistemas Distribuidos}
\UASetDate{2026}


\UASetClassNumber{04}
\UASetClassTitle{Consenso distribuido}
\UASetDocKind{Guía resuelta}

\title{Clase 04 --- Guía resuelta}
\author{\UAProfessor}
\date{\UADate}

\begin{document}

\section{Criterio de lectura}
Las resoluciones son modelos de razonamiento. Deben verificarse contra los supuestos del caso y no memorizarse como recetas independientes del dominio.

\section{Ejercicios y resoluciones completas}
\begin{uaexercise}
Definir formalmente el problema de consenso. Indicar claramente entradas, salidas y condiciones del sistema.
\end{uaexercise}
\begin{uaanswer}
La idempotencia separa intentos físicos de una operación lógica. Se usa una clave estable, se persiste el resultado junto con el efecto y los reintentos devuelven ese mismo resultado. El costo es estado adicional, política de expiración y coordinación con el almacenamiento.

La evidencia apropiada sería términos, votos, logs, commitIndex y un invariante de seguridad. También debe indicarse una condición que refute la elección o haga preferible una solución más simple.
\end{uaanswer}

\begin{uaexercise}
Explicar con precisión las propiedades:
Agreement, Validity y Termination.  
Dar un ejemplo concreto donde se viole cada una.
\end{uaexercise}
\begin{uaanswer}
La idempotencia separa intentos físicos de una operación lógica. Se usa una clave estable, se persiste el resultado junto con el efecto y los reintentos devuelven ese mismo resultado. El costo es estado adicional, política de expiración y coordinación con el almacenamiento.

La evidencia apropiada sería términos, votos, logs, commitIndex y un invariante de seguridad. También debe indicarse una condición que refute la elección o haga preferible una solución más simple.
\end{uaanswer}

\begin{uaexercise}
Explicar por qué consenso es necesario en sistemas replicados.
Relacionar con la clase anterior.
\end{uaexercise}
\begin{uaanswer}
La idempotencia separa intentos físicos de una operación lógica. Se usa una clave estable, se persiste el resultado junto con el efecto y los reintentos devuelven ese mismo resultado. El costo es estado adicional, política de expiración y coordinación con el almacenamiento.

El argumento debe identificar la hipótesis, construir la cadena lógica o el contraejemplo y concluir exactamente qué propiedad queda demostrada. Una intuición sin secuencia verificable no es suficiente.

La evidencia apropiada sería términos, votos, logs, commitIndex y un invariante de seguridad. También debe indicarse una condición que refute la elección o haga preferible una solución más simple.
\end{uaanswer}

\begin{uaexercise}
Dado un sistema con 3 nodos, construir una ejecución donde dos nodos decidan valores distintos.
¿Qué propiedad se viola?
\end{uaexercise}
\begin{uaanswer}
La idempotencia separa intentos físicos de una operación lógica. Se usa una clave estable, se persiste el resultado junto con el efecto y los reintentos devuelven ese mismo resultado. El costo es estado adicional, política de expiración y coordinación con el almacenamiento.

La evidencia apropiada sería términos, votos, logs, commitIndex y un invariante de seguridad. También debe indicarse una condición que refute la elección o haga preferible una solución más simple.
\end{uaanswer}

\begin{uaexercise}
Definir estado univalente y bivalente.
\end{uaexercise}
\begin{uaanswer}
La respuesta debe situarse en consenso, FLP, mayorías, Raft, Paxos, seguridad, progreso y máquina de estados replicada. Se comienza declarando el supuesto decisivo y el comportamiento observable; luego se recorre una ejecución nominal y otra bajo crash, demora arbitraria, elección dividida y pérdida temporal de mayoría. El mecanismo elegido debe vincularse con una garantía concreta, evitando afirmar propiedades fuera de su alcance.

La evidencia apropiada sería términos, votos, logs, commitIndex y un invariante de seguridad. También debe indicarse una condición que refute la elección o haga preferible una solución más simple.
\end{uaanswer}

\begin{uaexercise}
Explicar por qué el estado inicial puede ser bivalente.
\end{uaexercise}
\begin{uaanswer}
La respuesta debe situarse en consenso, FLP, mayorías, Raft, Paxos, seguridad, progreso y máquina de estados replicada. Se comienza declarando el supuesto decisivo y el comportamiento observable; luego se recorre una ejecución nominal y otra bajo crash, demora arbitraria, elección dividida y pérdida temporal de mayoría. El mecanismo elegido debe vincularse con una garantía concreta, evitando afirmar propiedades fuera de su alcance.

El argumento debe identificar la hipótesis, construir la cadena lógica o el contraejemplo y concluir exactamente qué propiedad queda demostrada. Una intuición sin secuencia verificable no es suficiente.

La evidencia apropiada sería términos, votos, logs, commitIndex y un invariante de seguridad. También debe indicarse una condición que refute la elección o haga preferible una solución más simple.
\end{uaanswer}

\begin{uaexercise}
Explicar el rol del scheduler adversario en la prueba FLP.
\end{uaexercise}
\begin{uaanswer}
FLP no afirma que consenso sea imposible en la práctica; demuestra que en asincronía total, con una posible falla y algoritmo determinista, no puede garantizarse terminación en todas las ejecuciones. Los sistemas reales recuperan progreso con sincronía parcial, timeouts y aleatoriedad, sin renunciar a seguridad.

La evidencia apropiada sería términos, votos, logs, commitIndex y un invariante de seguridad. También debe indicarse una condición que refute la elección o haga preferible una solución más simple.
\end{uaanswer}

\begin{uaexercise}
Construir una ejecución donde el sistema nunca decide.
\end{uaexercise}
\begin{uaanswer}
La respuesta debe situarse en consenso, FLP, mayorías, Raft, Paxos, seguridad, progreso y máquina de estados replicada. Se comienza declarando el supuesto decisivo y el comportamiento observable; luego se recorre una ejecución nominal y otra bajo crash, demora arbitraria, elección dividida y pérdida temporal de mayoría. El mecanismo elegido debe vincularse con una garantía concreta, evitando afirmar propiedades fuera de su alcance.

La evidencia apropiada sería términos, votos, logs, commitIndex y un invariante de seguridad. También debe indicarse una condición que refute la elección o haga preferible una solución más simple.
\end{uaanswer}

\begin{uaexercise}
Explicar por qué agregar timeouts rompe el modelo de FLP.
\end{uaexercise}
\begin{uaanswer}
La idempotencia separa intentos físicos de una operación lógica. Se usa una clave estable, se persiste el resultado junto con el efecto y los reintentos devuelven ese mismo resultado. El costo es estado adicional, política de expiración y coordinación con el almacenamiento.

El argumento debe identificar la hipótesis, construir la cadena lógica o el contraejemplo y concluir exactamente qué propiedad queda demostrada. Una intuición sin secuencia verificable no es suficiente.

La evidencia apropiada sería términos, votos, logs, commitIndex y un invariante de seguridad. También debe indicarse una condición que refute la elección o haga preferible una solución más simple.
\end{uaanswer}

\begin{uaexercise}
Describir detalladamente el proceso de elección de líder en Raft.
\end{uaexercise}
\begin{uaanswer}
La idempotencia separa intentos físicos de una operación lógica. Se usa una clave estable, se persiste el resultado junto con el efecto y los reintentos devuelven ese mismo resultado. El costo es estado adicional, política de expiración y coordinación con el almacenamiento.

La evidencia apropiada sería términos, votos, logs, commitIndex y un invariante de seguridad. También debe indicarse una condición que refute la elección o haga preferible una solución más simple.
\end{uaanswer}

\begin{uaexercise}
Explicar el rol de los términos (terms).
\end{uaexercise}
\begin{uaanswer}
La idempotencia separa intentos físicos de una operación lógica. Se usa una clave estable, se persiste el resultado junto con el efecto y los reintentos devuelven ese mismo resultado. El costo es estado adicional, política de expiración y coordinación con el almacenamiento.

La evidencia apropiada sería términos, votos, logs, commitIndex y un invariante de seguridad. También debe indicarse una condición que refute la elección o haga preferible una solución más simple.
\end{uaanswer}

\begin{uaexercise}
Describir el protocolo AppendEntries.
\end{uaexercise}
\begin{uaanswer}
La respuesta debe situarse en consenso, FLP, mayorías, Raft, Paxos, seguridad, progreso y máquina de estados replicada. Se comienza declarando el supuesto decisivo y el comportamiento observable; luego se recorre una ejecución nominal y otra bajo crash, demora arbitraria, elección dividida y pérdida temporal de mayoría. El mecanismo elegido debe vincularse con una garantía concreta, evitando afirmar propiedades fuera de su alcance.

La evidencia apropiada sería términos, votos, logs, commitIndex y un invariante de seguridad. También debe indicarse una condición que refute la elección o haga preferible una solución más simple.
\end{uaanswer}

\begin{uaexercise}
Definir formalmente la propiedad de Log Matching.
\end{uaexercise}
\begin{uaanswer}
La idempotencia separa intentos físicos de una operación lógica. Se usa una clave estable, se persiste el resultado junto con el efecto y los reintentos devuelven ese mismo resultado. El costo es estado adicional, política de expiración y coordinación con el almacenamiento.

La evidencia apropiada sería términos, votos, logs, commitIndex y un invariante de seguridad. También debe indicarse una condición que refute la elección o haga preferible una solución más simple.
\end{uaanswer}

\begin{uaexercise}
Explicar por qué el commit requiere mayoría.
\end{uaexercise}
\begin{uaanswer}
La respuesta debe situarse en consenso, FLP, mayorías, Raft, Paxos, seguridad, progreso y máquina de estados replicada. Se comienza declarando el supuesto decisivo y el comportamiento observable; luego se recorre una ejecución nominal y otra bajo crash, demora arbitraria, elección dividida y pérdida temporal de mayoría. El mecanismo elegido debe vincularse con una garantía concreta, evitando afirmar propiedades fuera de su alcance.

El argumento debe identificar la hipótesis, construir la cadena lógica o el contraejemplo y concluir exactamente qué propiedad queda demostrada. Una intuición sin secuencia verificable no es suficiente.

La evidencia apropiada sería términos, votos, logs, commitIndex y un invariante de seguridad. También debe indicarse una condición que refute la elección o haga preferible una solución más simple.
\end{uaanswer}

\begin{uaexercise}
Construir una ejecución donde un líder cae y otro lo reemplaza.
\end{uaexercise}
\begin{uaanswer}
La idempotencia separa intentos físicos de una operación lógica. Se usa una clave estable, se persiste el resultado junto con el efecto y los reintentos devuelven ese mismo resultado. El costo es estado adicional, política de expiración y coordinación con el almacenamiento.

La evidencia apropiada sería términos, votos, logs, commitIndex y un invariante de seguridad. También debe indicarse una condición que refute la elección o haga preferible una solución más simple.
\end{uaanswer}

\begin{uaexercise}
Describir los roles de Paxos: proposer, acceptor, learner.
\end{uaexercise}
\begin{uaanswer}
La respuesta debe situarse en consenso, FLP, mayorías, Raft, Paxos, seguridad, progreso y máquina de estados replicada. Se comienza declarando el supuesto decisivo y el comportamiento observable; luego se recorre una ejecución nominal y otra bajo crash, demora arbitraria, elección dividida y pérdida temporal de mayoría. El mecanismo elegido debe vincularse con una garantía concreta, evitando afirmar propiedades fuera de su alcance.

La evidencia apropiada sería términos, votos, logs, commitIndex y un invariante de seguridad. También debe indicarse una condición que refute la elección o haga preferible una solución más simple.
\end{uaanswer}

\begin{uaexercise}
Explicar la fase Prepare.
\end{uaexercise}
\begin{uaanswer}
La respuesta debe situarse en consenso, FLP, mayorías, Raft, Paxos, seguridad, progreso y máquina de estados replicada. Se comienza declarando el supuesto decisivo y el comportamiento observable; luego se recorre una ejecución nominal y otra bajo crash, demora arbitraria, elección dividida y pérdida temporal de mayoría. El mecanismo elegido debe vincularse con una garantía concreta, evitando afirmar propiedades fuera de su alcance.

La evidencia apropiada sería términos, votos, logs, commitIndex y un invariante de seguridad. También debe indicarse una condición que refute la elección o haga preferible una solución más simple.
\end{uaanswer}

\begin{uaexercise}
Explicar la fase Accept.
\end{uaexercise}
\begin{uaanswer}
La respuesta debe situarse en consenso, FLP, mayorías, Raft, Paxos, seguridad, progreso y máquina de estados replicada. Se comienza declarando el supuesto decisivo y el comportamiento observable; luego se recorre una ejecución nominal y otra bajo crash, demora arbitraria, elección dividida y pérdida temporal de mayoría. El mecanismo elegido debe vincularse con una garantía concreta, evitando afirmar propiedades fuera de su alcance.

La evidencia apropiada sería términos, votos, logs, commitIndex y un invariante de seguridad. También debe indicarse una condición que refute la elección o haga preferible una solución más simple.
\end{uaanswer}

\begin{uaexercise}
Definir formalmente qué significa que un valor sea elegido.
\end{uaexercise}
\begin{uaanswer}
La respuesta debe situarse en consenso, FLP, mayorías, Raft, Paxos, seguridad, progreso y máquina de estados replicada. Se comienza declarando el supuesto decisivo y el comportamiento observable; luego se recorre una ejecución nominal y otra bajo crash, demora arbitraria, elección dividida y pérdida temporal de mayoría. El mecanismo elegido debe vincularse con una garantía concreta, evitando afirmar propiedades fuera de su alcance.

La evidencia apropiada sería términos, votos, logs, commitIndex y un invariante de seguridad. También debe indicarse una condición que refute la elección o haga preferible una solución más simple.
\end{uaanswer}

\begin{uaexercise}
Explicar por qué dos valores distintos no pueden ser elegidos.
\end{uaexercise}
\begin{uaanswer}
La respuesta debe situarse en consenso, FLP, mayorías, Raft, Paxos, seguridad, progreso y máquina de estados replicada. Se comienza declarando el supuesto decisivo y el comportamiento observable; luego se recorre una ejecución nominal y otra bajo crash, demora arbitraria, elección dividida y pérdida temporal de mayoría. El mecanismo elegido debe vincularse con una garantía concreta, evitando afirmar propiedades fuera de su alcance.

El argumento debe identificar la hipótesis, construir la cadena lógica o el contraejemplo y concluir exactamente qué propiedad queda demostrada. Una intuición sin secuencia verificable no es suficiente.

La evidencia apropiada sería términos, votos, logs, commitIndex y un invariante de seguridad. También debe indicarse una condición que refute la elección o haga preferible una solución más simple.
\end{uaanswer}

\end{document}
